\begin{enumerate}[itemsep=1em]
	\item L'aire du carré est : $\mathcal{A}_\text{carré} = c \times c = 8~\text{cm} \times 8~\text{cm} = 64~\text{cm}^2$.
	\item Pour obtenir l'aire grisée, il faut soustraire l'aire de deux demi-disques de 8~cm de diamètre (donc l'aire d'un disque de 4~cm de rayon) à l'aire du carré : $\mathcal{A}_\text{grisée} = \mathcal{A}_\text{carré} - \pi\times R^2 = (64 - 16\pi)~\text{cm}^2$.
	
	\medskip
	\textit{($\pi$ est un nombre irrationnel, il n'a donc pas d'écriture décimale exacte. Ici, tout résultat décimal serait approché et donc faux dans cette question qui attendait un résultat exact.)}
	\item 
	\begin{enumerate}
		\item À l'échelle 125 : 1,
		 
		 1~cm est représenté par 125~cm donc \\
		 8~cm sont représentés par : $8\times125~\text{cm}=1~000~\text{cm}=10~\text{m}$.
		 
		 \medskip
		Le carré aura un côté de 10~m.
		\item On nomme $ABCD$ ce carré de 10~m de côté. Le triangle $ABC$ est donc rectangle en $B$ (car un carré a 4 angles droits) et d'après le théorème de Pythagore, on a :
		
		$AC^2=AB^2+BC^2$\\
		$AC^2=(10~\text{m})^2+(10~\text{m})^2$\\
		$AC^2=200~\text{m}^2$\\
		$AC=\sqrt{200}~\text{m}$\\
		$AC\approx14,14~\text{m}$ à 1 cm près
		
		\medskip
		La diagonale de ce carré a une longueur d'environ 14,14~m.
	\end{enumerate}
	
	\item %D'après la question 2, l'aire grisée du petit carré est : $(64 - 8\pi)~\text{cm}^2$. La figure dans la cour est un agrandissement de rapport 125, donc toutes les longueurs sont multipliées par 125 et les aires par 125\up{2}.
	
	Avec le même raisonnement qu'à la question 2, on a : 
	
	$\mathcal{A}_\text{grisée} = 10~\text{m} \times 10~\text{m} - \pi\times (5~\text{m})^2 \approx 21,4~\text{m}^2$.
	
	\medskip
	Il faut deux couches de peinture, il faudra donc peindre environ 42,8~m\up{2}.
	
	\medskip
	$\dfrac{42,8~\text{m}^2}{7~\text{m\up{2}/L}}\approx6,1~\text{L}$
	
	\medskip
	Il faudra utiliser environ 6,1~L de peinture.
	
	\medskip
	$6,1~\text{L}\div750~\text{mL}=6,1~\text{L}\div0,75~\text{L}\approx8,1$
	
	\medskip
	Finalement, il faudra prévoir 9 pots de peinture.
	
	\end{enumerate}